\section{Introduction}
\label{sec:intro}

Large simulations run as thousands of processes that share no memory and communicate by
sending messages. Almost all of them use the Message Passing Interface
(MPI)~\citep{mpi50}. MPI programs are usually deterministic by construction: every
receive names the process it expects a message from. One feature breaks this. A receive
may pass the wildcard \anysrc and accept the first matching message from any sender. Work
pools use it to serve whichever worker finishes first. Solvers use it when the set of
neighbours is not known in advance (AMG), and reduction trees use it to collect partial
sums in arrival order (Nekbone). Wildcard receives appear in \nProxyWild of
\nProxy US exascale proxy applications~\citep{sultana21proxy}.

\Cref{fig:tworuns} shows how little it takes for a wildcard to matter. Rank~0 receives
once from anyone and then once from rank~2. Ranks~1 and~2 each send one message. In
run~A, rank~1's message arrives first: the wildcard takes it, the second receive takes
rank~2's message, and the program finishes. In run~B, rank~2's message arrives first and
the wildcard takes it. The receive from rank~2 then waits forever, and rank~1's message
is never received. Both runs are legal MPI. Which one happens is decided by the network,
the load on the node and the MPI library's progress engine. The bad order may appear once
in thousands of test runs, disappear when the test is repeated under a debugger, and
then surface at scale. Documented cases include a race in a laser--plasma code that
appeared only at 524\,288 processes and took months to diagnose~\citep{laguna15cacm},
\pending{and the other §5.2 cases with their primary sources}.

\begin{figure}[t]
% Figure 1: the program of idiom 04 in two runs.
\centering
\begin{tikzpicture}[x=1.15cm,y=0.62cm,font=\footnotesize]
  % ---- run A
  \begin{scope}
    \node[anchor=west,font=\small\bfseries] at (-0.9,1.15) {Run A: finishes};
    \foreach \r/\y in {0/0,1/-1,2/-2} {
      \node[anchor=east] at (-0.15,\y) {rank \r};
      \draw[rank] (0,\y) -- (3.6,\y);
    }
    \draw[msg] (0.5,-1) -- (1.3,-0.08);
    \draw[msg] (1.4,-2) -- (2.5,-0.08);
    \node[wild] at (1.35,0) {};
    \node[ev] at (2.55,0) {};
    \node[above,font=\scriptsize\ttfamily] at (1.35,0.12) {Recv(ANY)};
    \node[above,font=\scriptsize\ttfamily] at (2.65,0.12) {Recv(2)};
  \end{scope}
  % ---- run B
  \begin{scope}[xshift=5.4cm]
    \node[anchor=west,font=\small\bfseries] at (-0.9,1.15) {Run B: hangs};
    \foreach \r/\y in {0/0,1/-1,2/-2} {
      \node[anchor=east] at (-0.15,\y) {rank \r};
      \draw[rank] (0,\y) -- (3.6,\y);
    }
    \draw[msg] (0.4,-2) -- (1.3,-0.08);
    \node[wild] at (1.35,0) {};
    \draw[alt] (1.2,-1) -- (2.3,-0.25);
    \node[font=\scriptsize,text=pendingcol,anchor=west] at (2.35,-0.62) {never received};
    \draw[pendingcol, line width=3pt] (2.0,0) -- (3.6,0);
    \node[above,font=\scriptsize\ttfamily] at (1.35,0.12) {Recv(ANY)};
    \node[above,font=\scriptsize\ttfamily,text=pendingcol] at (2.8,0.12) {Recv(2) waits};
  \end{scope}
\end{tikzpicture}

\caption{One program, two runs. Rank~0 posts \mpicall{Recv(ANY)} and then
\mpicall{Recv(from 2)}. Lines are ranks with time to the right; arrows are messages; the
ring marks the wildcard receive. In run~B the wildcard takes rank~2's message, the second
receive never completes and rank~1's message is lost.}
\label{fig:tworuns}
\end{figure}

Existing tools cover this poorly. Runtime checkers such as \MUST~\citep{hilbrich12must}
watch the run that happened. When that run takes the good order, there is nothing for them to
report. Exhaustive verifiers such as ISP~\citep{vakkalanka08isp},
DAMPI~\citep{vo10dampi}, MPI-SV~\citep{yu20mpisv} and CIVL~\citep{siegel15civl}
enumerate the schedules of a program. They pay for the number of competing
wildcards and stay at small scales. A line of predictive analyses starting with
MOPPER~\citep{forejt14mopper,forejt17mopper} takes one trace and asks whether some other
matching deadlocks~\citep{huang20ase,huang23tpds,li24parco}. These analyses target deadlocks
only and encode each query for a solver. They also assume that the program's sequence of
calls does not depend on what it receives. We show that this assumption makes them
unsound on ordinary code, such as a master that replies to whoever sent the last result.

Shared-memory concurrency faced the same gap and closed it with \emph{sound predictive
analysis}. From one trace of a multithreaded program, an analysis proves that a
reordering of the trace is a run of every program consistent with
it~\citep{kini17wcp,mathur18shb,mathur21syncp,tunc23spd,ang24strong}. It decides this in
linear time for well-chosen classes of reorderings~\citep{mathur20lics,kulkarni21concur}.
The soundness of these analyses now has a common
definition~\citep{liu25soundness}. That machinery does not transfer verbatim. Its
central object, the reads-from relation, becomes a \emph{matching} of receives to sends.
A matching is consumed exactly once and constrained by MPI's non-overtaking rule. Whether
a send completes depends on the buffering threshold of the MPI implementation. And after
a receive takes a different message, the receiver's control flow may change.

\paragraph{This paper.} We define which alternative matchings one trace certifies, prove
that the certification is sound for every program consistent with the trace, and build
\tool, a tool that predicts such matchings, classifies their consequences, and confirms
them by replay. The analysis keeps a causally closed prefix of the trace, the
\emph{horizon}. Everything in it happens identically under the alternative, and nothing
after a changed receive is assumed. Inside the horizon an alternative is feasible iff a
constraint graph over issue, completion and match events is acyclic. Each prediction
carries the set of buffering thresholds at which it holds. Each finding states its level
of evidence. \Lone findings follow from the trace alone. \LoneSP findings additionally
assume the single-path continuation of earlier work. \Ltwo findings use a skeleton of
the program's dependence on message metadata.

\paragraph{Contributions.}
\begin{itemize}[leftmargin=*]
\item A model of programs consistent with a trace and a constraint graph with match nodes.
  Its acyclicity over the horizon certifies an alternative matching for every such
  program (\cref{thm:sound}). For single-change alternatives, the check is also complete
  relative to the reorderings that keep every other match (\cref{thm:complete}).
\item The buffering profile. Feasibility of a single-change alternative is monotone in
  the buffering threshold above the threshold at which the observed run is realisable
  (\cref{thm:mono}). Deadlock predictions are neither upward- nor downward-closed
  (\cref{prop:nonmono}), so a prediction must name its thresholds.
\item Levels of evidence. Single-path prediction is unsound on code whose control flow
  depends on message metadata (\cref{lem:sp}). Deadlocks contained in the horizon are
  predicted without that assumption (\cref{thm:l1dead}).
\item \tool: a PMPI tracer, an analyser whose check is quasi-linear in the trace per
  alternative (\cref{thm:cost}), and a replay library. The replay forces a
  predicted matching, logs that the forced match happened, and is checked against a
  control run.
\item An evaluation on the MPI Bugs Initiative and \MBB against \MUST and against repeated
  randomised runs, a differential check of soundness against an exhaustive explorer, and
  two real applications. \pending{RQ results after the rerun.}
\end{itemize}
